% The original template (from Trevor) had a custom \appendix command,
% but I found it to break figure/table counters. I'm not sure how
% reliable my fix is, so I ended up reverting back to the standard
% latex version, and renaming the custom command to \myappendix.  You
% can try both and see how things work out:
% 1) Call \appendix once, and then make each appendix a \chapter
% 2) Call \myappendix once, and then make each appendix a \section.

\appendix
\chapter{}
\label{appendix:A}
% Appendix sections, subsections etc should not appear in the table of 
% contents. So set tocdepth to 0 and set it back to 2 at the end. Do 
% this for each appendix.
%\addtocontents{toc}{\protect\setcounter{tocdepth}{0}}

\newcommand{\maintitle}[1]{%
        \noindent{\Large\textbf{#1}}
        \vspace{0.2cm}
        \hrule
        \vspace{0.5cm}
}

\maintitle{PROJECT PROPOSAL}


\section{Description of Project}
%Quantum Error Correction is a field required to bridge the hardware implementation of quantum computers to the error rates required for practical Fault Tolerant Quantum Computing (FTQC). Color Codes are an example of a QEC code that encodes a logical qubit as a collection of physical qubits.
Trapped Ion (TI) systems are a possible platform for quantum systems, of which Quantum Charge-Coupled Devices (QCCDs) are a modular architectural implementation which has demonstrated potential for Quantum Error Correction (QEC) on small TI systems \cite{ryananderson2021realizationrealtimefaulttolerantquantum}. Recently, a tooflow was proposed which provides a framework in which QEC codes can be compiled onto different choices of QCCD hardware to simulate and test metrics such as logical error rates \cite{scottPaper}. While this tool provides rudimentary support for codes such as the Surface Code, compilation of Color Codes are still a niche and under-researched area. This project will focus on building a set of compiler passes for this toolbox for evaluating Color Codes as a candidate QEC code on different QCCD hardware choices.

\subsection{Main Goals}
The main goal of this project is to create a compiler that can take in QCCD hardware parameters to produce a representation of a color code that can be used to run QEC using primitive QCCD operations.
\newline
\noindent{The tools/adaptations that are required are:}
\begin{itemize}
    \item A compiler pass to map logical color code qubits to QCCD ions
    \item A compiler pass to translate color code circuit operations to QCCD operations
    \item Adaptation and addition to the simulation and test infrastructure to run the compiled color code circuit
\end{itemize}

This project will be considered successful if a user can pass in a choice of QCCD hardware parameters and produce estimate logical error rates for a given color code.

\subsubsection{Key Challenges}
The main challenge of this task is to be able to leverage features of a color code to provide a suitable and ideally near-optimal mapping onto different QCCD architecture.

\subsection{Extensions}
A few stretch goals for this include some research and further tool extensions:

\begin{itemize}
    \item Include support for different Color code tessellations e.g. [4,8,8]
    \item Provide mappings for similar topological codes
    \item Investigate the implementation and ideal QCCD hardware choices relating to higher dimensional QEC codes
    \item Investigate ideal QCCD hardware choices for FTQC using color codes
    \item Experimenting with different qubit mapping and ion routing approaches
\end{itemize}

\section{Evaluation}\label{eval}
Evaluation of the project will be done as follows:
\begin{itemize}
    \item The tool for mapping the color code will be compared against how close the simulated QEC cycles gets to the optimal mapping for a pre-determined set of individual color codes and QCCD device models.
    \item The compiler will be evaluated with a set of QEC cycles and tested for correct termination as well as the correct sequences of gates applied to each ion from a set of color codes and QCCD device models.
    \item Measurement of compilation time and the number of routing operations required with comparison to theoretic times
    \item Evaluation of logical error rates for color code, QCCD hardware configuration pairs
\end{itemize}

%\noindent{Further stretch goal(s) evaluation:}

%\begin{itemize}
%    \item The effect of QCCD architectural features/choices on elapsed/QEC round time, logical error rates and similar relevant metrics
%    \item The fault-tolerance threshold of color code and QCCD configuration pairs
%    \item Effect of mapping algorithm to compilation time and number of routing operations required
%\end{itemize}
\section{Success Criteria}
\begin{enumerate}
    \item A compilation/mapping scheme for color codes onto QCCD architecture that can be used to estimate logical error rates.
    \item Suitable simulation and testing infrastructure to carry out the evaluation  on cycles of Color Codes on QCCD devices and the compiler.
\end{enumerate}


\section{Risk Management}
As quantum computing is a complex field, one source of risk is to lack knowledge in the theory to understand or implement a part of this project. To mitigate this, I will ensure to regularly read up on literature to reinforce my knowledge and consult experts such as lecturers and my supervisor when facing any confusion.

\section{Resource Declaration}
\begin{itemize}
    \item \textbf{Computing resources:}
        \begin{itemize}
            \item ThinkPad T430 2-core CPU with 8GB RAM (Windows 10 and Linux) for portability
            \item Custom PC with a 6-core CPU and 32GB RAM (Windows 11)
        \end{itemize}
    \item \textbf{Dissertation Backup:} Majority of my dissertation will be written and saved on \href{https://www.overleaf.com}{Overleaf}. I will also save the project onto an external hard-disk at intermediate stages.

    \item \textbf{Code Backup:} Code will be stored in a remote GitHub repository. Local development will be done by branching from the main code, writing code, then merged back into main.
\end{itemize}

\section{Starting Point}
Over the summer, I completed a Coursera on QEC by Google Quantum AI along with some exploration of the stim python library and the \href{https://eur03.safelinks.protection.outlook.com/?url=https%3A%2F%2Fgithub.com%2Fscottjones03%2Fpublic-material-2025%2Ftree%2Fmain&data=05%7C02%7Cme526%40cam.ac.uk%7C4332f74bac054c3f6e0908ddf3b5f0cd%7C49a50445bdfa4b79ade3547b4f3986e9%7C1%7C0%7C638934684941570154%7CUnknown%7CTWFpbGZsb3d8eyJFbXB0eU1hcGkiOnRydWUsIlYiOiIwLjAuMDAwMCIsIlAiOiJXaW4zMiIsIkFOIjoiTWFpbCIsIldUIjoyfQ%3D%3D%7C0%7C%7C%7C&sdata=qcgGlEcD0yaOa%2BshDJWb3NB59bPOL3gSbxWLu0kjf%2FM%3D&reserved=0}{QEC test framework repository} on GitHub.

\section{Timetable}
\begin{enumerate}
    \item \textbf{10/10 - 23/10: Preparation}
    
    Deliverables:
    \begin{itemize}
        \item A comprehensive report on color codes including motivation and requirements for mapping color codes to QCCD
        \item Tests run on the QEC testing framework to demonstrate sufficient understanding of the current implementation
    \end{itemize}

    \item \textbf{24/10 - 6/11: Mapping Research}
    
    Deliverables: Research into methods of mapping qubits to ions
    \begin{itemize}
        \item Research and selection of mapping methodologies such as clustering, greedy heuristics and satisfiability (SMT)
        \item Report on paper readings 
        \item Draft of preparation chapter
    \end{itemize}
    The algorithm should make use of the characteristics of the color code.
    
    \item \textbf{7/11 - 4/12: Mapping Implementation}
    
    Deliverables: mappings of color codes to the available QCCD hardware choices. 
    \begin{itemize}
        \item Assignment of color code qubits to physical qubits in hardware
        \item Adjustment and addition to framework to compile color code circuits
        \item Writing testing and evaluation infrastructure
        \item Draft of the Implementation chapter
    \end{itemize}
    Extended timeframe to account for 3 module assignments that take place over mid-late November.
    
    \item \textbf{5/11 - 25/12: Extension Work}
    Deliverables: Implementation of any relevant extension work into core and discussion of results in relevant chapters.
    
    \item \textbf{26/12 - 15/1:  Testing + Evaluation}
    Deliverables: evaluating implementations and generating graphs and metrics. Draft of the evaluation chapter
   
    Some parameters that may be tested for the QCCD QEC implementation with enough time:
    \begin{itemize}
        \item Communication topology
        \item Segment lengths
        \item Trap sizings
        \item Color Code tessellations
        \item Code dimensions        
        \item Ancilla to data qubit ratio
        \item Distinct QCCD memory and processing zones
    \end{itemize}

    \item \textbf{16/1 - 29/1: Dissertation Intro}

    Deliverable: Draft of introduction chapter. Draft of progress report due for 6/2. Send whole draft dissertation to supervisor for checks.

    \item \textbf{30/1 - 19/2: Editing and Conclusion Chapter}
    
    Deliverable: Draft of conclusion chapter. Finalising and reviewing previously drafted chapters. Progress report submission and presentation.

    Slack added to account for high weekly lecture load in Lent term.
    
    \item \textbf{20/2 - 26/3: Finalising + Editing of Chapters}

    Deliverable: Full dissertation draft sent to DoS and Supervisor before the beginning of Easter break.
    
    \item \textbf{27/3 - 15/5: Refine Final Dissertation and Submission}
\end{enumerate}

%\addtocontents{toc}{\protect\setcounter{tocdepth}{2}}

%%% Local Variables: ***
%%% mode: latex ***
%%% TeX-master: "thesis.tex" ***
%%% End: ***
